% !TeX root = ../navier-stokes-manuscript.tex
\subsection{The critical-height problem inside a stretching vortex}\label{sub:ThePerfectlyStretchingVortex}

Take a region of fluid carrying coherent rotation and place it in an
extensional strain.  Write
\[
  \omega:=\nabla\times u,
  \qquad
  S:=\frac12\bigl(\nabla u+(\nabla u)^{\mathsf T}\bigr).
\]
The curl of the momentum equation gives
\begin{equation}\label{eq:BodyVorticityEquation}
  \partial_t\omega+(u\cdot\nabla)\omega
  =S\omega+\nu\Delta\omega.
\end{equation}
The term \(S\omega\) is the exact action of strain on vorticity.  When
\(\omega\) points along an extensional direction of \(S\),
\[
  \omega\cdot S\omega>0,
\]
and the rotation intensifies.

The vortex does not receive that amplification without changing its shape.
Incompressibility says
\[
  \operatorname{tr}S=\nabla\cdot u=0.
\]
Extension along the vortex axis must be accompanied by contraction in the
transverse directions.  The structure becomes longer, its cross-section
becomes narrower, and the velocity variation carried across that
cross-section is pressed into less space.  If \(\Delta U\) is the variation
across a width \(\ell\), then
\[
  |\nabla u|\sim\frac{\Delta U}{\ell}.
\]
The regularity problem is already visible in this ratio.  The amount of
velocity may stay finite while the distance over which it changes collapses.

Kinetic energy is allowed to miss that event.  Its exact balance is
\begin{equation}\label{eq:BodyEnergyIdentity}
  \frac12\norm{u(t)}_2^2
  +\nu\int_0^t\norm{\nabla u(s)}_2^2\,ds
  =\frac12\norm{u_0}_2^2.
\end{equation}
The first term measures velocity over volume.  It does not retain the width
across which that velocity changes.  A narrowing region can carry less and
less kinetic energy while its gradients become steeper.

Viscosity cannot be discussed apart from that narrowing.  Pair
\eqref{eq:BodyVorticityEquation} with \(\omega\).  The pointwise form is
\[
  \frac12(\partial_t+u\cdot\nabla)|\omega|^2
  =\omega\cdot S\omega
  +\frac\nu2\Delta|\omega|^2
  -\nu|\nabla\omega|^2.
\]
After integration over \(\R^3\), transport and the Laplacian flux disappear,
and the competition becomes
\begin{equation}\label{eq:BodyEnstrophyBalance}
  \frac12\frac d{dt}\norm{\omega(t)}_2^2
  +\nu\norm{\nabla\omega(t)}_2^2
  =\int_{\R^3}\omega\cdot S\omega\,dx.
\end{equation}
Stretching and diffusion are acting on the same rotating structure.  As its
width falls, both terms strengthen.  The balance has shown us the mechanism,
but it has not told us whether viscosity gains on the stretching.

\divider{Shrink the entire equation}

Now shrink the picture without changing the law that governs it.  For
\(\lambda>0\), define
\begin{equation}\label{eq:BodyNavierStokesScaling}
  u_\lambda(x,t):=\lambda u(\lambda x,\lambda^2t),
  \qquad
  p_\lambda(x,t):=\lambda^2p(\lambda x,\lambda^2t).
\end{equation}
Every term in the momentum equation receives the factor \(\lambda^3\):
\[
\begin{aligned}
  \partial_tu_\lambda
  &=\lambda^3(\partial_tu)(\lambda x,\lambda^2t),
  &
  (u_\lambda\cdot\nabla)u_\lambda
  &=\lambda^3((u\cdot\nabla)u)(\lambda x,\lambda^2t),
  \\
  \nabla p_\lambda
  &=\lambda^3(\nabla p)(\lambda x,\lambda^2t),
  &
  \Delta u_\lambda
  &=\lambda^3(\Delta u)(\lambda x,\lambda^2t).
\end{aligned}
\]
Nothing in the smaller copy has handed viscosity an extra power of scale.
Transport, pressure, and diffusion remain locked in the original proportion.
The scale change has instead shortened the clock by \(\lambda^{-2}\).

To see what the shrinking leaves unchanged, let
\(\Lambda=(-\Delta)^{1/2}\).  A change of variables gives
\begin{equation}\label{eq:BodyHomogeneousSobolevScaling}
  \norm{u_\lambda(t)}_{\dot H^s}^2
  =\lambda^{2s-1}\norm{u(\lambda^2t)}_{\dot H^s}^2.
\end{equation}
The exponent vanishes at one height:
\[
  2s-1=0,
  \qquad
  s=\frac12.
\]
We call
\begin{equation}\label{eq:BodyCriticalHeightDefinition}
  H(t):=\frac12\norm{u(t)}_{\dot H^{1/2}}^2
  =\frac12\norm{\Lambda^{1/2}u(t)}_2^2
\end{equation}
the critical height.  It is the spatial reading that a lawful concentration
cannot make smaller by shrinking itself.  Energy sits below it and loses one
power of \(\lambda\); the first derivative sits above it and gains one:
\[
  \norm{u_\lambda(t)}_2^2
  =\lambda^{-1}\norm{u(\lambda^2t)}_2^2,
  \qquad
  H_\lambda(t)=H(\lambda^2t),
  \qquad
  \norm{\nabla u_\lambda(t)}_2^2
  =\lambda\norm{\nabla u(\lambda^2t)}_2^2.
\]
This is why finite energy can remain calm while a derivative sharpens.  The
critical height occupies the point between those two readings where the
concentration is still fully visible.

One rescaled profile makes the separation concrete.  For a nonzero
divergence-free \(\phi\in\Ssigma\), set
\[
  u_\ell(x):=\ell^{-1}\phi(x/\ell).
\]
Then
\[
\begin{aligned}
  \norm{u_\ell}_\infty&\sim\ell^{-1},
  &\norm{\nabla u_\ell}_\infty&\sim\ell^{-2},
  \\
  \norm{u_\ell}_2^2&\sim\ell,
  &\norm{u_\ell}_{\dot H^{1/2}}^2
  &=\norm{\phi}_{\dot H^{1/2}}^2.
\end{aligned}
\]
The amplitude and gradient rise, the energy falls, and the critical height
does not move.  This family is not a Navier--Stokes trajectory.  It shows the
shape of the loophole left open by energy and respected by scaling.

\divider{What the equation does at critical height}

Scaling has selected the right height.  The equation now tells us how that
height moves.  For
\[
  E_s(t):=\frac12\norm{\Lambda^su(t)}_2^2,
\]
keep transport and pressure together in the signed response
\[
  \mathcal P_s(t)
  :=-
  \operatorname{Re}
  \pair{\Lambda^su}
  {\Lambda^s\bigl((u\cdot\nabla)u+\nabla p\bigr)}_{L^2}.
\]
The pressure part vanishes in this global pairing after incompressibility is
used.  We keep it visible because the scalar balance came from the complete
momentum equation, not from transport alone.  Pairing at height \(s\) gives
\begin{equation}\label{eq:BodySpatialEnergyBalanceIntro}
  E_s'=\mathcal P_s-2\nu E_{s+1}.
\end{equation}
At the critical height,
\begin{equation}\label{eq:BodyCriticalHeightBalanceIntro}
  H'=\mathcal P_{1/2}-2\nu E_{3/2}.
\end{equation}
Positive production raises \(H\) only after paying the viscous loss beside
it.  The equation has placed the unresolved stretching--diffusion contest at
the scale-invariant height where neither side can escape by rescaling.

\divider{Give the shrinking scale a clock}

The linear heat flow assigns an exact time scale to a width.  If
\(v_t=\nu\Delta v\), then
\[
  \widehat v(\xi,t+\delta t)
  =e^{-\nu|\xi|^2\delta t}\widehat v(\xi,t).
\]
A structure of width \(r\) lives near \(|\xi|\simeq r^{-1}\).  Diffusion
changes its amplitude by order one on the time
\begin{equation}\label{eq:BodyHeatComparisonTime}
  \tau_\nu(r)\asymp\frac{r^2}{\nu}.
\end{equation}
This is not yet an estimate on the nonlinear fluid.  It is the amount of time
viscosity assigns to a scale once the evolution has actually formed that
scale.

Take geometrically shrinking widths
\[
  r_n:=2^{-n}r_0,
  \qquad
  \tau_n:=\frac{r_n^2}{\nu}.
\]
Their heat clocks satisfy
\begin{equation}\label{eq:BodyZenoHeatClock}
  \tau_n=\frac{r_0^2}{\nu}4^{-n},
  \qquad
  \sum_{n=0}^{\infty}\tau_n
  =\frac{4r_0^2}{3\nu}<\infty.
\end{equation}
This is the Zeno opening: the clocks attached to infinitely many smaller
scales fit inside a finite amount of elapsed time.

The arithmetic does not create a singularity.  The fluid would have to form
\(r_0\), then \(r_1\), then \(r_2\), continuing without losing the
stretching alignment while viscosity is already acting on every scale that
has appeared.  Each new event would begin from the velocity produced by the
previous one.  The widths shrink, the clocks shorten, and the critical
stretching--diffusion balance has to be reproduced again.

That is what turns the spatial picture into a question about response.  If the
velocity is to cross sharper states during shorter intervals, what must its
acceleration do?  If the acceleration is forced upward, what must lead it?
The singularity has brought us to the temporal derivative Tower.
